\section{Conclusion}

In this experiment, several important phenomena related to the polarization of light were successfully investigated. Malus’s law was 
experimentally verified, showing the expected cosine-squared dependence between transmitted intensity and the angle between polarizer 
and analyzer. Although small deviations from the theoretical prediction occurred, mainly due to non-ideal optical components and 
experimental uncertainties, the overall agreement confirmed the validity of the model.

The birefringence of calcite crystals demonstrated how anisotropic materials split light into ordinary and extraordinary rays with 
different polarization states. Furthermore, the behavior of quarter-wave plates illustrated the conversion between linear, circular, 
and elliptical polarization, depending on the orientation of the optical axis.

Using the Jones formalism, polarization states and optical elements could be described mathematically, providing a useful theoretical 
framework for understanding the observed effects. The experiment with sucrose solutions showed that circular birefringence can be applied 
to determine concentrations through optical rotation measurements.

Finally, the Faraday effect in glass was successfully observed. A linear relationship between the rotation angle and the magnetic field 
strength was found for all investigated wavelength ranges. The calculated Verdet constants decreased with increasing wavelength, which 
agrees with the theoretical expectations for dispersive optical media.

Overall, the experiment provided a comprehensive understanding of polarization phenomena and demonstrated the close connection between 
theoretical optics and experimental observations.
